\documentclass[11pt]{article}
\usepackage[a4paper,margin=25mm]{geometry}
\usepackage{amsmath,amssymb,amsthm,mathtools}
\usepackage[colorlinks=true,urlcolor=blue]{hyperref}
\newtheorem{theorem}{Theorem}
\newtheorem{definition}[theorem]{Definition}
\newtheorem{remark}[theorem]{Remark}
\newcommand{\Z}{\mathbb Z}
\newcommand{\B}{\mathbb B}
\newcommand{\Sp}{\operatorname{Spec}}
\title{The Original Third Element: Exact Arithmetic, Support, and Spectral Receivers}
\author{}
\date{19 September 2026}
\begin{document}
\maketitle

\paragraph{Origin and scope.}
The original idea was to adjoin what its originator called a ``trivial number'' to the numbers and investigate its operations. The July 2025 precursor \cite{original} uses an adjoined $\mathcal T$, with addition leaving the other argument unchanged and multiplication returning $\mathcal T$. The sequential-adjunction manuscript \cite{sequential}, Section 9 in its source \texttt{5.tex}, explicitly adds $\Omega$ absorbing both operations. Section 10 records the three singlets $\{\mathcal T,0_{\Z},\Omega\}$. The definitions, rather than the subsequent term ``Split-Zero'', identify this history. We use $\mathcal T$ throughout this note to retain the original notation.

The original $\Omega$ extension is checked here in precisely its stated algebraic category. No global law $0x=0$ is inserted. All maps below specify their identities and algebraic signatures. The calculations establish structural content and an exact comparison with a known absorbing-element construction; they assert neither historical novelty nor an analytic estimate for the RH programme.

\begin{definition}
Let $\mathcal W$ be the category whose objects $(A,+,\cdot,0_A,1_A)$ have a commutative additive monoid, a commutative multiplicative monoid, and distributivity. Morphisms preserve both binary operations and both specified identities. The equation $0_Aa=0_A$ is not required. This is the explicitly specified version of the source's ``commutative unital hemiring'' convention.
For an object $A$, let $A^\Omega=A\sqcup\{\Omega_A\}$, with old operations on $A$ and with any sum or product involving $\Omega_A$ equal to $\Omega_A$.
\end{definition}

Fix a nonzero commutative unital ring $R$, and set
\[
S_R=R\sqcup\{\mathcal T\},\qquad H_R=S_R^\Omega,
\qquad e=0_R\in R.
\]
On $S_R$ use the original ring operations for two arguments in $R$, and use
\[
\mathcal T+x=x+\mathcal T=x,\qquad
\mathcal T x=x\mathcal T=\mathcal T\quad(x\in S_R).
\]
Thus the three elements $\mathcal T,e,\Omega$ are distinct in $H_R$.

\begin{theorem}[The original extension retains the original algebra]
Both $S_R$ and $H_R$ are objects of $\mathcal W$, and the inclusion $S_R\hookrightarrow H_R$ is an injective morphism in $\mathcal W$. The identities of $H_R$ are $\mathcal T$ and $1_R$. In particular,
\[
\mathcal T\Omega=\Omega,\quad e\Omega=\Omega,
\quad \mathcal T+\Omega=\Omega,
\]
while $\mathcal T x=\mathcal T$ for every $x\in S_R$.
\end{theorem}
\begin{proof}
The map $S_R\to R\times\B$ given by $r\mapsto(r,1)$ and $\mathcal T\mapsto(0,0)$ identifies it with the subset $(R\times\{1\})\cup\{(0,0)\}$. This subset contains the two identities and is closed under componentwise sum and product: two supported terms give $(r+s,1)$ and $(rs,1)$, and $(0,0)$ has the required identity/absorber rules. Hence all the standard semiring axioms, and therefore the $\mathcal W$ axioms, hold on $S_R$.

For any $A\in\mathcal W$, associativity and commutativity on $A^\Omega$ hold on triples in $A$ by inheritance. A triple containing $\Omega_A$ evaluates to $\Omega_A$ under either association for each operation. For distributivity $x(y+z)=xy+xz$, all three arguments in $A$ give the old law. If $x=\Omega_A$, both sides are $\Omega_A$; if $y$ or $z$ is $\Omega_A$, the left side is $\Omega_A$ and the right side has a summand $\Omega_A$, so is also $\Omega_A$. The old additive identity still fixes $\Omega_A$ under addition, and the old multiplicative identity still fixes it under multiplication. These exhaust the cases. Apply this to $A=S_R$. Restriction to $S_R$ leaves every operation and both identities unchanged, proving the inclusion assertion and the displayed formulas.
\end{proof}

\begin{theorem}[Two exact observers and complete reconstruction]
Write $R^\Omega=R\sqcup\{\bot\}$ for the absorbing extension of the ring, and $C=\B^\Omega=\{0,1,\infty\}$. The maps
\[
\begin{array}{c|ccc}
x&r\in R&\mathcal T&\Omega\\\hline
q(x)&r&0_R&\bot\\
\sigma(x)&1&0&\infty
\end{array}
\]
are surjective morphisms $q:H_R\to R^\Omega$ and $\sigma:H_R\to C$ in $\mathcal W$. The pair $(q,\sigma)$ is an isomorphism onto the subobject
\[
E_R=\{(r,1):r\in R\}\ \cup\ \{(0_R,0),(\bot,\infty)\}
\subset R^\Omega\times C.
\]
The only nonsingleton fibre of $q$ is $\{\mathcal T,e\}$. Consequently
\[
H_R/(\mathcal T\sim e)\simeq R^\Omega.
\]
The fibre $\sigma^{-1}(1)=R$ recovers the entire original ring, with its original zero $e$, unit, operations, and additive inverses.
\end{theorem}
\begin{proof}
For two inputs in $R$, $q$ preserves operations by the ring laws. For one input $\mathcal T$ and the other $r\in R$, addition gives $r$ on both sides and multiplication maps $\mathcal T r=\mathcal T$ to $0_R=0_Rr$. For two inputs $\mathcal T$, both identities give $0_R$. When either input is $\Omega$, the target's $\bot$ absorbs both operations, exactly as required. The identities are preserved.

For $\sigma$, two inputs in $R$ have sum and product in $R$, even when their ordinary amplitude is zero; the target computations are $1+1=1$ and $1\cdot1=1$. An input $\mathcal T$ uses Boolean zero as additive identity and multiplicative absorber against $0$ or $1$. An input $\Omega$ uses the global absorber $\infty$, including $0\infty=\infty$. Thus $\sigma$ preserves both operations and identities. Both maps are visibly surjective.

The displayed three types of pair are disjoint. The first coordinate distinguishes all pairs of the first type, so $(q,\sigma)$ is injective, and its image is exactly $E_R$. This proves its isomorphism assertion including closure and preservation of identities. The fibre description of $q$ follows directly from the table. Its kernel equivalence is a congruence identifying only $\mathcal T$ with $e$, hence the indicated quotient is exact. Finally the fibre of $1$ is literally $R$ with its unchanged operations; the restricted ring zero is $e$, which is distinguished from the whole object's additive identity $\mathcal T$.
\end{proof}

\paragraph{Known absorbing extension.}
Bergstra and Tucker \cite[Section 2.1]{BT} define $\mathsf{Enl}_{\bot}(R)$ by adjoining to a commutative ring an element absorbing all its operations. On the $(+,\cdot,0,1)$ signature, $R^\Omega$ above is exactly that construction under $\bot\mapsto\bot$ and the identity on $R$. If negation is included, set $-\bot=\bot$. Thus $q$ is an exact quotient onto this cited construction. Division belongs to additional structure in common meadows \cite{BP}; no division is being inferred here.

\begin{theorem}[What standard linear observers can receive]
If $D$ is a standard commutative unital semiring and $f:H_R\to D$ is a $\mathcal W$ morphism, then $D$ is the one-element semiring. The unique map $H_R\to\{0=1\}$ is therefore its universal map to standard semirings.
\end{theorem}
\begin{proof}
Since $f(\mathcal T)=0_D$ and $\Omega\mathcal T=\Omega$, one has
\[
f(\Omega)=f(\Omega)f(\mathcal T)=f(\Omega)0_D=0_D.
\]
Applying $f$ to $\Omega+1_R=\Omega$ then gives $0_D+1_D=0_D$, so $1_D=0_D$. Every $d\in D$ satisfies $d=d1_D=d0_D=0_D$. The singleton is itself a standard semiring and the unique map to it preserves all operations and identities, establishing the universal assertion. The previous theorem exhibits nontrivial receivers in $\mathcal W$ and an injective joint receiver, so this result specifies exactly which extra axiom causes the collapse.
\end{proof}

\begin{definition}[The source's ideal convention]
An ideal is a nonempty subset closed under addition and under multiplication by every element of the ambient object. A prime ideal is a proper ideal $P$ such that $ab\in P$ implies $a\in P$ or $b\in P$. In $\mathcal W$ this convention does not additionally require the additive identity to belong to the ideal. Use the closed sets $V(I)=\{P:I\subseteq P\}$ for the resulting prime spectrum. These are the conventions used in the original sequential manuscript's spectral calculation.
\end{definition}

\begin{theorem}[Exact retained spectrum and the extra generic point]
The ideals of $S_R$ are $\{\mathcal T\}$ and $J\cup\{\mathcal T\}$, for ring ideals $J\subseteq R$. Its primes are $\{\mathcal T\}$ and $\mathfrak p\cup\{\mathcal T\}$, for $\mathfrak p\in\Sp R$.

The ideals of $H_R$ are $\{\Omega\}$ and $I\cup\{\Omega\}$, for ideals $I$ of $S_R$. Its primes are consequently
\[
\{\Omega\},\qquad \{\mathcal T,\Omega\},\qquad
\mathfrak p\cup\{\mathcal T,\Omega\}\quad(\mathfrak p\in\Sp R).
\]
The map $P\mapsto P\cup\{\Omega\}$ is a homeomorphism $\Sp S_R\to V_{H_R}(\mathcal T)$. The remaining point $\{\Omega\}$ is open and dense. Thus, for finite Krull dimension,
\[
\dim S_R=\dim R+1,\qquad \dim H_R=\dim R+2.
\]
In particular for $R=\Z$ and a rational prime $p$ the full strict chain is
\[
\{\Omega\}\subsetneq\{\mathcal T,\Omega\}
\subsetneq\{0_{\Z},\mathcal T,\Omega\}
\subsetneq p\Z\cup\{\mathcal T,\Omega\},
\]
and the dimension is $3$.
\end{theorem}
\begin{proof}
Any ideal $I$ of $S_R$ contains $\mathcal T$, by multiplying one of its elements by $\mathcal T$. If $I\cap R$ is empty then $I=\{\mathcal T\}$. Otherwise $J=I\cap R$ is closed under ring addition and multiplication by $R$, and under additive inverses since $-1\in R$. It is therefore a ring ideal. Conversely every such $J\cup\{\mathcal T\}$ is an ideal by the operation definitions. A product of two elements of $R$ stays in $R$, so $\{\mathcal T\}$ is prime. For $J\cup\{\mathcal T\}$, primality reduces exactly to primality of $J$ in $R$.

Every ideal $L$ of $H_R$ contains $\Omega$, by multiplying an element of $L$ by $\Omega$. If $L\cap S_R$ is nonempty, it is an ideal of $S_R$ because the old operations are unchanged; if it is empty, $L=\{\Omega\}$. Conversely the asserted sets are all ideals by direct closure. A product equals $\Omega$ precisely when a factor equals $\Omega$, so $\{\Omega\}$ is prime. For the other ideals, pairs of arguments in $S_R$ test exactly the old primality condition, while a pair containing $\Omega$ already has a factor in the ideal. This proves the prime classification.

Every prime except $\{\Omega\}$ contains $\mathcal T$, giving the displayed image $V_{H_R}(\mathcal T)$. On that image, closed sets correspond under the ideal extension $I\mapsto I\cup\{\Omega\}$ and contraction to $S_R$, by the already established ideal classification. This proves the homeomorphism. The remaining point is $D_{H_R}(\mathcal T)$, so is open. It is dense because its ideal is contained in every prime, and hence its closure is the whole spectrum. The prime poset is obtained from $\Sp R$ by prepending the two displayed strict minimum points. Every chain arises this way or is a subchain. Prefixing those two points to a maximal-length ring-prime chain gives equality of the dimensions. For $\Z$ the ring-prime chains have maximal length one, and the displayed chain explicitly realizes length three.
\end{proof}

\begin{remark}[A precise limitation on contracting primes]
For the inclusion $S_R\hookrightarrow H_R$, contraction of the new prime $\{\Omega\}$ is empty, so it does not define a prime of $S_R$ under the stated nonempty-ideal convention. The proved homeomorphism on $V_{H_R}(\mathcal T)$ is the exact spectral receiver. No total contravariant map on these spectra is inferred from the inclusion.
\end{remark}

\paragraph{Assessment.}
Renaming $\Omega$ as $\infty$ and retaining the entire operation table produces an isomorphic object by the map fixing $S_R$ and sending $\Omega$ to $\infty$. Its content is the specified adjunction. The exact additional state survives in the three-valued support observer, the paired arithmetic/support reconstruction and the extra spectral point under the source's ideal convention. Identifying $\mathcal T$ with $e$ has a precisely described quotient onto the established absorbing extension of a ring. These are verified ways to retain and compare the original idea; no inference from the symbol $\infty$ supplies a topology, order, or analytic bound.

\begin{thebibliography}{9}
\bibitem{original} HI-AI, \emph{An Algebraic Structure Incorporating a Z/1Z-Symmetric Element Adjoined to the Integers: Construction, Analysis, and Generalizations}, recorded publication date 12 July 2025, Zenodo version \href{https://zenodo.org/records/17555345}{17555345}, source \texttt{11.tex}, Section 3. The concept and original name are supplied by the project's originator; the record labels the manuscript a creative exploration.
\bibitem{sequential} HI-AI, \emph{An Examination of Semiring Structures Derived from the Integers by Sequential Adjunction of Identity and Absorbing Elements}, Zenodo version \href{https://zenodo.org/records/17547186}{17547186}, source \texttt{5.tex}, Sections 9--10 and 12. The version contains multiple source files; all claims attributed here identify \texttt{5.tex} specifically.
\bibitem{BT} Jan A. Bergstra and John V. Tucker, \emph{On the Axioms of Common Meadows: Fracterm Calculus, Flattening and Incompleteness}, The Computer Journal 66(7) (2023), 1565--1572, Section 2.1, \href{https://doi.org/10.1093/comjnl/bxac026}{doi:10.1093/comjnl/bxac026}.
\bibitem{BP} Jan A. Bergstra and Alban Ponse, \emph{Division by Zero in Common Meadows}, \href{https://arxiv.org/abs/1406.6878}{arXiv:1406.6878}, version 3. Cited for the broader common-meadow context, not as a proof of the split-support receivers established above.
\end{thebibliography}
\end{document}
